% Chapter questions for sections 07_modular, 08_modexp, 09_gcd, 09b_ext_euclid (end-sem split). Numbers verified by script (scratchpad/r4).
\section{Modular Arithmetic, Exponentiation, GCD and Extended Euclid}\label{q:07_modular}
\noindent{\small Notes: \S\ref{sec:modular}, \S\ref{sec:modexp}, \S\ref{sec:gcd}, \S\ref{sec:exteuclid}.}

\subsection*{Check your understanding}
\begin{enumerate}[leftmargin=1.6em,itemsep=3pt]
  \item \textbf{Differentiate $a \equiv b \pmod n$ and $a \bmod n = b$.} The congruence means $n \mid (a - b)$, i.e.\ $a$ and $b$ leave the same remainder, and either side may be any member of the residue class, so $38 \equiv 14 \pmod{12}$ is true. The equation with the operator says that $b$ is \emph{the} least non-negative remainder, so $38 \bmod 12 = 2$ and ``$38 \bmod 12 = 14$'' is false.
  \item \textbf{Find $-23 \bmod 7$ and explain the correction step.} Truncated division gives $-23 = (-3) \times 7 + (-2)$, whose remainder is negative. Decreasing $q$ by 1 and adding 7 to $r$ leaves the value unchanged: $-23 = (-4) \times 7 + 5$, so $-23 \bmod 7 = \mathbf{5}$.
  \item \textbf{Why is congruence modulo $n$ an equivalence relation, and what are its classes?} It is reflexive ($n \mid 0$), symmetric ($n \mid a - b \Rightarrow n \mid b - a$) and transitive ($a - c = (a - b) + (b - c)$). It therefore splits $\mathbb{Z}$ into $n$ disjoint residue classes $[0], [1], \dots, [n-1]$, where $[a] = \{a + kn : k \in \mathbb{Z}\}$; $\mathbb{Z}_n$ takes the least non-negative member of each class.
  \item \textbf{Prove $(a \times b) \bmod n = [(a \bmod n)(b \bmod n)] \bmod n$ and say why cryptography needs it.} With $a = q_1 n + r_1$ and $b = q_2 n + r_2$, $ab = (q_1 q_2 n + q_1 r_2 + q_2 r_1)n + r_1 r_2$, which differs from $r_1 r_2$ by a multiple of $n$. It allows reduction after every multiplication, so in RSA and Diffie--Hellman no intermediate value ever exceeds $n^2$.
  \item \textbf{Why can a congruence not simply be divided by a common factor?} Because $6 \equiv 16 \pmod{10}$, yet dividing by 2 gives $3 \not\equiv 8 \pmod{10}$. Division means multiplying by an inverse, so cancelling a factor $c$ is valid only when $\gcd(c, n) = 1$.
  \item \textbf{Differentiate additive and multiplicative inverses in $\mathbb{Z}_n$.} The additive inverse $b$ satisfies $a + b \equiv 0$; it always exists ($b = n - a$) and gives subtraction, e.g.\ decrypting a shift cipher. The multiplicative inverse satisfies $ab \equiv 1$; it exists only when $\gcd(a, n) = 1$ and gives division, e.g.\ in the affine, Hill and RSA ciphers. Hence $\mathbb{Z}_n$ is used for additive inverses and $\mathbb{Z}_n^*$ for multiplicative ones.
  \item \textbf{Show that 6 has no multiplicative inverse modulo 15.} $\gcd(6, 15) = 3$. If $6x \equiv 1$, then $6x - 15k = 1$ for some $k$, but the left side is always a multiple of 3, so this is impossible.
  \item \textbf{How many solutions does $6x \equiv 9 \pmod{15}$ have in $\mathbb{Z}_{15}$? Find them.} $d = \gcd(6, 15) = 3$ divides 9, so there are 3. Dividing by 3 gives $2x \equiv 3 \pmod 5$; $2^{-1} \bmod 5 = 3$, so $x_0 = 9 \bmod 5 = 4$, and the solutions are $x_0 + 5k$: $\mathbf{4, 9, 14}$ (check: $6 \times 9 = 54 = 3 \times 15 + 9$).
  \item \textbf{Why may the base of $a^b \bmod m$ be reduced modulo $m$ but not the exponent?} The base is a factor, and the product rule lets any factor be replaced by a congruent one. The exponent only counts factors: $2^5 \equiv 2 \pmod 3$ but $2^{5 \bmod 3} = 4 \equiv 1$. When $\gcd(a, m) = 1$ the exponent may be reduced modulo $\phi(m)$ instead (Euler).
  \item \textbf{How many modular multiplications does square-and-multiply need for $a^{45}$, and why is the method polynomial-time?} $45 = 101101_2$ has 6 bits and four 1-bits: 5 squarings and 3 multiplications, \textbf{8} in all against 44 (the left-to-right column $c$ runs $1, 2, 5, 11, 22, 45$). In general it needs at most $2\lfloor\log_2 b\rfloor$ multiplications, linear in the bit length $k$ of the exponent, while repeated multiplication needs $b - 1 \approx 2^k$.
  \item \textbf{Prove that $\gcd(a, b) = \gcd(b, a \bmod b)$.} Write $a = qb + r$. A common divisor of $a$ and $b$ divides $a - qb = r$, and a common divisor of $b$ and $r$ divides $qb + r = a$. The two pairs have the same common divisors, hence the same greatest one.
  \item \textbf{Why does Euclid's algorithm terminate, and why is it fast?} The second number is replaced by a remainder smaller than itself, so the sequence of remainders strictly decreases to 0. Every two steps the remainder at least halves (if $r_{i+1} > r_i/2$ the quotient is 1 and $r_{i+2} = r_i - r_{i+1} < r_i/2$), so at most about $2\log_2 b$ divisions are needed.
  \item \textbf{State the invariant of the extended Euclidean algorithm and explain why the final $t_1$ is $b^{-1} \bmod n$.} In every row $r_1 = s_1 n + t_1 b$ and $r_2 = s_2 n + t_2 b$, because $r$, $s$, $t$ are updated by the same rule (new $=$ first $- q \times$ second). At the end $r_1 = \gcd(n, b)$; if it is 1, then $s_1 n + t_1 b = 1$, and reducing modulo $n$ gives $t_1 b \equiv 1$.
  \item \textbf{Find $17^{-1} \bmod 60$ with the table.} Quotients $3, 1, 1, 8$; remainders $9, 8, 1, 0$; $t$ values $-3, 4, -7, 60$. The final $r_1 = 1$ and $t_1 = -7$, so $17^{-1} = -7 + 60 = \mathbf{53}$ (check: $17 \times 53 = 901 = 15 \times 60 + 1$).
  \item \textbf{When does $ax + by = c$ have integer solutions, and why infinitely many?} Exactly when $d = \gcd(a, b)$ divides $c$, since $d$ divides every $ax + by$. From one solution $(x_0, y_0)$, the pair $(x_0 + k\,b/d,\ y_0 - k\,a/d)$ is again a solution for every integer $k$, because the two changes cancel in $ax + by$.
\end{enumerate}
\noindent{\small Longer practice for these chapters: Practice Problems sets I, J, K and P (\S\ref{sec:practice}) and drill types 4--6 (\S\ref{sec:drill}).}
